\documentclass{article}
\usepackage{geometry}[left=1cm,right=1cm]
\usepackage{eso-pic,graphicx} % Required for inserting images
\usepackage[default]{frcursive}

\title{Inhabitants of Farlorret\\
\large A report of the creatures surrounding Farlorret}
\date{\today}

\AddToShipoutPictureBG{\includegraphics[width=\paperwidth,height=\paperheight]{../../Y7.png}}


\begin{document}
    \maketitle

    \section*{Western lower fields}

    Here rather calm ecosystems thrive, within the shade of the forests small critters of a wide variety can be found taking shelter from the larger predators roaming the skies. 

    Said predators pertain mostly to the scavenging Death swans. Feeding mostly on corpses of larger beasts but is not afraid to produce their own meals if needed. 

    \section*{Western elemental plateau}

    Here the veil between the elemental planes of earth and wind grows thin. Small pebble like elementals litter the plateau with larger ones collecting towards their titan up by the northern mountains.

    Other scavengers are found here specifically Hyenas are known to roam in larger packs with leaders which grows to the size of a small horse.

    \section*{Eastern forests}

    The eastern forests contain similar ecology as the western if it was not for the recent destruction of a large swat of habitat in a fire. Directly south lives goblins operating both a ferry and an elevator leading north.
    
    Meanwhile, the southern part of have a lower terrain and increased hydration which makes the areas more prone to a more swampish biome with lizard-kin who live there.


    \section*{Far eastern drake burrows}

    Beyond the forest a large water source is frequented by drakes and its kin. Some of these creatures seems to have more or less of a connection to the elemental planes which makes the classification between elementals and beasts difficult.

    South towards the sea the terrain drops and makes for a valley where sightings of rat-kin who tame these drakes. 

    \section*{Northern plateau}

    Here lives dryads, harpies and survivors of different kinds in the ruins of the great ruins. The forest is lively and bountiful with predators as large and exotic as owl bears prowling beneath its boughs.

    Sparse, unconfirmed reports further north describes orange clay or sand beyond the mountains. This would indicate a mesa or other dry habitats perhaps conducive to scouting for further habitats.

    \section*{The sea}

    The sea is bountiful and deep. This makes its habitants hard to survey, crustacians of many kinds do show up on the shores and archipelagos. It is somewhat safe to speculate a multitude of different fishes and other marine lives hides beneath the waves.
    \end{document}